\section*{Chapitre \ref{ch:b3:lithosphere-thermal-models} --- Modèles thermiques de la lithosphère}

\begin{solution}{exo:b3:lithosphere-thermal-models:1}
$L^2/\kappa = \qty{e12}{s}$, \qty{e14}{s} et \qty{e16}{s}~: 0.032, 3.2 et
\qty{317}{Ma}. Un kilomètre oublie sa température en dizaines de milliers de
ans, une centaine de kilomètres en centaines de millions.
\end{solution}

\begin{solution}{exo:b3:lithosphere-thermal-models:2}
$u = \qty{e-3}{m}/\qty{3.156e7}{s} = \qty{3.17e-11}{m.s^{-1}}$, donc
$\mathrm{Pe} = 3.17\times10^{-11}\times\num{3e4}/10^{-6} = 0.95$. L'advection et
la conduction ont la même importance~: la roche qui remonte transporte la chaleur vers le haut
et raidit le géotherme près de la surface, la conduction seule
aurait tendance à sous-estimer le gradient.
\end{solution}

\begin{solution}{exo:b3:lithosphere-thermal-models:3}
$\sqrt{\kappa t} = \sqrt{\num{0.8e-6}\times5\times\num{3.156e13}} =
\qty{11236}{m}$, $\eta = 10000/(2\times11236) = 0.445$, et $T = 1450\times0.471
= \qty{683}{\celsius}$.
\end{solution}

\begin{solution}{exo:b3:lithosphere-thermal-models:4}
$490/\sqrt1 = 490$, $490/5 = 98$ et $490/10 = \qty{49}{mW.m^{-2}}$. La sonde
sous-estimerait très probablement la première valeur~: la croûte jeune est fissurée et ouverte à
l'eau de mer en circulation, qui emporte une partie de la chaleur par des sources chaudes.
\end{solution}

\begin{solution}{exo:b3:lithosphere-thermal-models:5}
$z = 2\times1.163\sqrt{\kappa t}$~: avec $\sqrt{\kappa t} = 15.89$, 31.78 et
\qty{50.25}{km}, $z = 37.0$, 73.9 et \qty{116.9}{km}. La couche refroidie double
de 10 à \qty{40}{Ma}~; dès \qty{100}{Ma}, elle est déjà plus profonde
que la plaque de \qty{95}{km}.
\end{solution}

\begin{solution}{exo:b3:lithosphere-thermal-models:6}
$C$ est proportionnel à $\alpha$~: $\alpha = \num{3.1e-5}\times350/370.5 =
\qty{2.93e-5}{K^{-1}}$. $C$ est aussi proportionnel à $T_m - T_0$ (un manteau à
$1450\times350/370.5 = \qty{1370}{\celsius}$) et à $\sqrt\kappa$ (une
diffusivité $0.8\times(350/370.5)^2 = \qty{0.71e-6}{m^2.s^{-1}}$)~: la profondeur seule
ne permet pas de les séparer, ce qui explique qu'on utilise aussi le flux de chaleur.
\end{solution}

\begin{solution}{exo:b3:lithosphere-thermal-models:7}
$3.1\times1450/\num{95e3} = \qty{47.3}{mW.m^{-2}}$ et $3.1\times1350/\num{125e3}
= \qty{33.5}{mW.m^{-2}}$. Le demi-espace les atteint à $(490/47.3)^2 =
\qty{107}{Ma}$ et $(490/33.5)^2 = \qty{214}{Ma}$~: la plaque mince et chaude conserve
le flux de chaleur du plancher ancien, à \qty{47}{mW.m^{-2}}, à partir d'environ \qty{100}{Ma},
tandis que la plaque plus épaisse et plus froide laisserait le flux tomber sous
\qty{40}{mW.m^{-2}}, valeur du demi-espace à \qty{150}{Ma}.
\end{solution}

\begin{solution}{exo:b3:lithosphere-thermal-models:8}
$A_0h_r = 0.4\times60 = \qty{24}{mW.m^{-2}}$, donc $A_0 =
\num{24e-3}/\num{e4} = \qty{2.4}{\micro W.m^{-3}}$ et $q_m = 60 - 24 =
\qty{36}{mW.m^{-2}}$. À \qty{35}{km}, $T = \num{36e-3}\times\num{35e3}/2.5 +
\num{2.4e-6}\times\num{e8}/2.5\times(1 - e^{-3.5}) = 504 + 96\times0.970 =
\qty{597}{\celsius}$.
\end{solution}

\begin{solution}{exo:b3:lithosphere-thermal-models:9}
$h_r = \Delta q_0/\Delta A_0 = \num{18e-3}/\num{1.8e-6} = \qty{10}{km}$, et
$q_m = q_0 - h_rA_0 = 66 - 30 = \qty{36}{mW.m^{-2}}$, puisque $h_rA_0 =
\num{e4}\times\qty{3.0e-6}{W.m^{-3}} = \qty{30}{mW.m^{-2}}$.
\end{solution}

\begin{solution}{exo:b3:lithosphere-thermal-models:10}
(a) $Q(t) = \int_0^tk(T_m - T_0)/\sqrt{\pi\kappa t'}\,\mathrm dt' =
2k(T_m - T_0)\sqrt{t/(\pi\kappa)}$. (b) La chaleur perdue est le contenu
extrait de la colonne, $Q = \rho_mc\int_0^\infty(T_m - T)\,\mathrm dz$, et
l'isostasie donne $(\rho_m - \rho_w)(d - d_r) = \rho_m\alpha\int_0^\infty(T_m -
T)\,\mathrm dz = \alpha Q/c$. (c) $c = 3.1/(3300\times\num{0.8e-6}) =
\qty{1174}{J.kg^{-1}.K^{-1}}$~; à \qty{40}{Ma}, $t = \qty{1.262e15}{s}$ et
$Q = 2\times3.1\times1450\times\sqrt{\num{1.262e15}/(\pi\times\num{0.8e-6})} =
\qty{2.015e14}{J.m^{-2}}$~; $d - d_r = \num{3.1e-5}\times\num{2.015e14}/
(1174\times2270) = \qty{2343}{m}$, soit $370.5\sqrt{40}$.
\end{solution}

\begin{solution}{exo:b3:lithosphere-thermal-models:11}
Le terme $n = 3$ vaut $3104\times(8/\pi^2)\times\tfrac19e^{-9t/36.2} =
280\,e^{-9t/36.2}$ mètres~: \qty{218}{m} à \qty{1}{Ma}, \qty{23}{m} à
\qty{10}{Ma} et \qty{1.9}{m} à \qty{20}{Ma}. Elle descend à \qty{10}{m} lorsque
$e^{-9t/36.2} = 10/280$, à $t = (36.2/9)\ln28 = \qty{13.4}{Ma}$.
\end{solution}

\begin{solution}{exo:b3:lithosphere-thermal-models:12}
(a) Pour $z \gg h_r$, $e^{-z/h_r} \to 0$. (b) $q_m = \qty{24}{mW.m^{-2}}$,
$A_0h_r^2/k = \num{16e-3}\times\num{e4}/2.5 = \qty{64}{K}$, et $z = (1300 -
64)\times2.5/\num{24e-3} = \qty{128.8}{km}$. (c) $q_m = 2.5\times1236/\num{250e3}
= \qty{12.4}{mW.m^{-2}}$~: le \omterm{def:b3:lithosphere-thermal-models:reduced}{flux de chaleur réduit} doit être à peu près divisé par deux, comme sous un
craton dont la racine épaisse et froide laisse peu de chaleur pénétrer par le bas.
\end{solution}

\begin{solution}{pb:b3:lithosphere-thermal-models:1}
\textbf{1.} $\sqrt{\num{0.8e-6}\times\num{3.156e14}} = \qty{15889}{m}$ et
$\sqrt{\num{0.8e-6}\times\num{1.2624e15}} = \qty{31778}{m}$.

\textbf{2.} $\eta = 20000/31778 = 0.6294$ et $20000/63556 = 0.3147$~:
$T = 1450\times0.6266 = \qty{909}{\celsius}$ à \qty{10}{Ma} et
$1450\times0.3437 = \qty{498}{\celsius}$ à \qty{40}{Ma}.

\textbf{3.} $z = 2\times0.7174\sqrt{\kappa t}$~: \qty{22.8}{km} et
\qty{45.6}{km}. Elle atteint \qty{60}{km} lorsque $\sqrt{\kappa t} = 60/1.4348 =
\qty{41.82}{km}$, à $t = (\num{41.82e3})^2/\num{0.8e-6} = \qty{2.186e15}{s}$,
\qty{69.3}{Ma}.

\textbf{4.} $q = 3.1\times1450/(\sqrt\pi\times15889) = \qty{159.6}{mW.m^{-2}}$ à
\qty{10}{Ma} et \qty{79.8}{mW.m^{-2}} à \qty{40}{Ma}~; $C_q = 159.6\sqrt{10}
= \qty{505}{mW.m^{-2}.Ma^{1/2}}$.

\textbf{5.} $u = \qty{7.92e-10}{m.s^{-1}}$, $\mathrm{Pe} = 7.92\times10^{-10}
\times\num{95e3}/\num{0.8e-6} = 94$. La chaleur se propage latéralement avec la plaque
cent fois plus vite qu'elle ne diffuse latéralement, si bien que chaque colonne
se refroidit uniquement par conduction verticale.

\textbf{6.} $(\num{95e3})^2/\num{0.8e-6} = \qty{1.13e16}{s}$, \qty{357}{Ma}.
À \qty{150}{Ma}, $\sqrt{\kappa t} = \qty{61.5}{km}$ et la couche refroidie, soit
plusieurs fois cette épaisseur, atteint bien en dessous de \qty{95}{km}~: un modèle de demi-espace appliqué à un plancher aussi ancien
suppose que le froid a atteint des profondeurs où le manteau convecte.

\textbf{7.} $(\rho_m - \rho_w)(d - d_r) = \rho_m\alpha\int(T_m - T)\,\mathrm dz$,
les masses de la colonne de dorsale et de la colonne ancienne étant égales au-dessus de
la profondeur de compensation. $T_m - T = T_m\operatorname{erfc}\eta$ tombe sous
\qty{0.5}{\percent} de $T_m$ pour $\eta > 2$, la partie profonde n'ajoute rien.

\textbf{8.} $C = 2\times3300\times\num{3.1e-5}\times1450/2270\times
\sqrt{\num{0.8e-6}\times\num{3.156e13}/\pi} = 0.13069\times2834.8 =
\qty{370.5}{m.Ma^{-1/2}}$~; $d = 2500 + 370.5\sqrt{10} = \qty{3672}{m}$ et
$2500 + 370.5\sqrt{40} = \qty{4843}{m}$.

\textbf{9.} Le modèle donne $2500 + 370.5\sqrt{42.5} = \qty{4915}{m}$,
\qty{440}{m} plus profond que la médiane. La médiane prend en compte la charge
sédimentaire, les monts sous-marins et les plateaux volcaniques, les renflements des points chauds, et
les cellules situées sur les flancs des zones de fracture et des marges, autant de facteurs qui rendent le plancher
moins profond.

\textbf{10.} $C$ est proportionnel à $T_m$~: $370.5\times1350/1450 =
\qty{344.9}{m.Ma^{-1/2}}$, proche de la valeur ajustée de 350.

\textbf{11.} L'eau de mer qui circule dans la croûte jeune fracturée transporte
une partie de la chaleur vers l'extérieur par des sources hydrothermales, entre les sondes.

\textbf{12.} $\tau = (\num{95e3})^2/(\pi^2\times\num{0.8e-6}) =
\qty{1.143e15}{s} = \qty{36.2}{Ma}$~; l'amplitude est $3300\times
\num{3.1e-5}\times1450\times\num{95e3}/4540 = \qty{3104}{m}$.

\textbf{13.} $e^{-40/36.2} = 0.33124$, $e^{-360/36.2}/9 = \num{5.3e-6}$ et le
terme $n = 5$ est encore plus petit~; le crochet vaut $1 - 0.81057\times0.33125 =
0.73150$, et $d = 2500 + 3104\times0.73150 = \qty{4771}{m}$.

\textbf{14.} Par le terme $n = 3$, $3104\times0.81057\times\num{5.3e-6} =
\qty{0.013}{m}$~: à \qty{40}{Ma}, $e^{-9t/\tau}$ est égal à $10^{-4}$ fois
$e^{-t/\tau}$, et seul le premier mode subsiste.

\textbf{15.} À \qty{40}{Ma}, $4843 - 4771 = \qty{73}{m}$~; à
\qty{50}{Ma}, $2500 + 370.5\sqrt{50} = 5120$ contre $2500 + 3104\times(1 -
0.81057\times0.25127) = 4972$~: \qty{148}{m}. Par interpolation, $40 + 10\times
(100 - 73)/(148 - 73) = \qty{43.6}{Ma}$, soit environ \qty{44}{Ma} (la série complète
donne 44.0).

\textbf{16.} $q = (3.1\times1450/\num{95e3})[1 + 2(0.33124 + 0.01204 +
0.00005)] = 47.32\times1.68666 = \qty{79.8}{mW.m^{-2}}$, identique à
celle du demi-espace~: à \qty{40}{Ma}, la base ne se fait pas encore sentir à la
surface.

\textbf{17.} $2500 + 3104 = \qty{5604}{m}$ et $kT_m/L = \qty{47.3}{mW.m^{-2}}$.

\textbf{18.} $d(120) = 2500 + aL\bigl[1 - (8/\pi^2)e^{-120/\tau}\bigr]$, avec
$a = \rho_m\alpha T_m/\bigl(2(\rho_m - \rho_w)\bigr) = \qty{32.673}{m/km}$ et
$\tau = L^2/(\pi^2\kappa)$.

\textbf{19.} $L = 3200/\bigl(32.673\,[1 - 0.81057e^{-120/\tau}]\bigr)$. À partir de
$L = \qty{100}{km}$~: $\tau = \qty{40.13}{Ma}$, crochet 0.95925, $L =
\qty{102.10}{km}$~; puis $\tau = 41.83$, $L = 102.67$~; $\tau = 42.30$, $L =
102.82$~; $\tau = 42.43$, $L = 102.87$~; il se stabilise à $L = \qty{102.9}{km}$.

\textbf{20.} $\tau = \qty{42.48}{Ma}$, $aL = \qty{3362}{m}$~: $d(40) = 2500 +
3362\times(1 - 0.81057e^{-40/42.48}) = \qty{4799}{m}$, contre \qty{4843}{m}
pour le demi-espace et une valeur mesurée de \qty{4800}{m}~: la plaque rend compte des deux âges.

\textbf{21.} $\tau = \qty{42.5}{Ma}$~; profondeur asymptotique $2500 + 3362 =
\qty{5862}{m}$.

\textbf{22.} $3362\times0.81057\,e^{-t/42.48} = 100$ donne $t = 42.48\ln27.25 =
\qty{140}{Ma}$.

\textbf{23.} La plaque qui convient au vieux plancher de l'océan fictif
est épaisse de \textbf{\qty[text-series-to-math=true]{102.9}{km}}.
\end{solution}
